%% ─────────────────────────────────────────────────────────────────────────────
%% CV — Talvin Ackbaraly (English Version, AI Engineer — Qantev)
%% Serif (XCharter), one restrained accent colour, thin rules, bullets, no icons
%% Compile: latexmk -pdf cv-talvin-ackbaraly-en.tex
%% ─────────────────────────────────────────────────────────────────────────────
\documentclass[10.5pt, a4paper]{extarticle}

\usepackage[T1]{fontenc}
\usepackage[utf8]{inputenc}
\usepackage[english]{babel}
\usepackage[top=1cm, bottom=1cm, left=1.6cm, right=1.6cm]{geometry}
\usepackage{XCharter}
\usepackage[letterspace=80]{microtype}
\usepackage{titlesec}
\usepackage{enumitem}
\usepackage{tabularx}
\usepackage{ragged2e}
\usepackage{xcolor}
\usepackage[scaled=0.94,noupquote]{inconsolata}
\usepackage{tikz}
\usepackage{varwidth}
\usepackage{graphicx}
\usepackage[hidelinks]{hyperref}

\hypersetup{pdftitle={Resume — Talvin Ackbaraly}, pdfauthor={Talvin Ackbaraly}}

%% ── Colour system: one accent, everything else neutral ───────────────────────
\definecolor{accent}{HTML}{2563A0}  % THE accent: name, sections, links, bullets
\definecolor{ink}{HTML}{1A1A1A}     % body text (warm near-black, not pure 000)
\definecolor{soft}{HTML}{6E6E6E}    % secondary information (context, tech)
\definecolor{hair}{HTML}{C9C9C9}    % section rules
\definecolor{chipbg}{HTML}{EAF0F7}  % tag pill background
\definecolor{chipfg}{HTML}{2F5B87}  % tag pill text

\hypersetup{colorlinks=true, urlcolor=accent, linkcolor=accent}
\AtBeginDocument{\color{ink}}

\pagestyle{empty}
\setlength{\parindent}{0pt}
\setlength{\parskip}{0pt}
\linespread{1.02}
\RaggedRight
\hyphenpenalty=10000
\exhyphenpenalty=10000
\emergencystretch=1em
\microtypesetup{expansion=false}

%% Section: letter-spaced small caps over a hairline
\titleformat{\section}
  {\large\scshape\lsstyle\color{accent}}{}{0pt}{}
  [\vspace{-5pt}{\color{hair}\rule{\textwidth}{0.4pt}}]
\titlespacing*{\section}{0pt}{13pt}{5pt}


%% Tag pill: \chip{CUDA}
\newcommand{\chip}[1]{%
  \tikz[baseline=(c.base)]{\node[inner xsep=4.5pt, inner ysep=2pt, rounded corners=2.5pt,
    fill=chipbg, text=chipfg, font=\footnotesize] (c) {#1};}\hspace{1pt}%
}
%% \entry{Title}{Date}{Subtitle}{Tags}
\newcommand{\entry}[4]{%
  \textbf{#1} \hfill #2\\
  \textit{#3} \hfill {#4}\par
}
%% \project{Name}{Context}{Tags}
\newcommand{\project}[3]{%
  \textbf{#1}{\color{soft}\enspace—\enspace\textit{#2}} \hfill {#3}\par
}
\newcommand{\gap}{\par\vspace{6pt}}


%% NB: at the head of a `p' cell, {\color{...}} shifts content down one line — always \textcolor.

%% ── 1. Education: stacked entries, no tags ───────────────────────────────
\newcommand{\edu}[4]{%
  \textbf{#1}\hfill #2\par\vspace{1.5pt}
  \textcolor{accent}{\textit{#3}}\par\vspace{1.5pt}
  {\small\color{soft}#4}\par
}

%% ── 2. Projects: prose (title first) / tags ──────────────────────────────────
%% Chaque ligne calcule sa propre largeur de colonne de tags : on mesure la
%% largeur naturelle des pastilles, on la plafonne, et la prose prend le reste.
\newlength{\projgapb}\setlength{\projgapb}{0.4cm}
\newlength{\projmaxtag}\setlength{\projmaxtag}{3.35cm}
\newlength{\projtagw}\newlength{\projprosew}
\newsavebox{\projtagbox}
%% #1 titre (en tête de prose, sans point final), #2 prose, #3 pastilles.
%% varwidth compose les pastilles jusqu'à \projmaxtag mais rétrécit à la largeur
%% réellement occupée : chaque ligne garde donc sa propre gouttière.
\newcommand{\proj}[3]{%
  \sbox\projtagbox{\begin{varwidth}[t]{\projmaxtag}\raggedleft #3\end{varwidth}}%
  \setlength{\projtagw}{\wd\projtagbox}%
  \setlength{\projprosew}{\dimexpr\textwidth-\projgapb-\projtagw\relax}%
  \noindent
  \parbox[t]{\projprosew}{\RaggedRight \textbf{#1.} #2}%
  \hspace{\projgapb}%
  \usebox\projtagbox
  \par
}
\newcommand{\projgap}{\vspace{10pt}}

%% ── 3. Experience: role first, full-width prose ───────────────────────
%% #1 role, #2 date, #3 employer, #4 pills, #5 prose.
%% The date takes the top-right corner — that is where the eye lands first;
%% the pills move down to the employer line.
\newcommand{\job}[5]{%
  \textbf{#1}\hfill #2\par\vspace{1.5pt}
  \textcolor{accent}{#3}\hfill {#4}\par\vspace{4pt}
  #5\par
}

%% ── 4. Skills: label + pills ─────────────────────────────────────
\newlength{\skilllab}\setlength{\skilllab}{4.95cm}
\newcommand{\skillrow}[2]{%
  \noindent\parbox[t]{\skilllab}{\textbf{#1}}%
  \parbox[t]{\dimexpr\textwidth-\skilllab\relax}{\RaggedRight #2}\par\vspace{3.5pt}
}


%% ── Header: round photo, top right ───────────────────────────────────
%% Photo: assets/photo.jpg; otherwise an empty placeholder is shown.
\newlength{\photosize}\setlength{\photosize}{3.4cm}
\IfFileExists{assets/photo.jpg}{\def\photofile{assets/photo.jpg}}{%
  \IfFileExists{../../assets/photo.jpg}{\def\photofile{../../assets/photo.jpg}}{}}
\newcommand{\headerphoto}{%
  \ifdefined\photofile
    \tikz{\clip (0,0) circle (0.5\photosize);
      \node at (0,0) {\includegraphics[width=\photosize]{\photofile}};
      \draw[hair, line width=0.6pt] (0,0) circle (0.5\photosize);}%
  \else
    \tikz{\draw[hair, fill=chipbg, line width=0.6pt] (0,0) circle (0.5\photosize);
      \node[text=chipfg, font=\footnotesize] at (0,0) {Photo};}%
  \fi
}

\begin{document}

\noindent
\begin{minipage}[c]{\dimexpr\textwidth-\photosize-0.3cm\relax}
  \RaggedRight
  {\color{accent}\fontsize{25}{30}\selectfont\scshape\textls[40]{Talvin Ackbaraly}}\\[6pt]
  \textbf{AI Engineer internship · 6 months from February 2027}\\[5pt]
  {\small
  \href{mailto:talvin.ackbaraly@icloud.com}{talvin.ackbaraly@icloud.com} \enspace·\enspace
  \href{tel:+33769389008}{+33 7 69 38 90 08} \enspace·\enspace
  \href{https://talvin-ackbaraly.com}{talvin-ackbaraly.com}\\[2pt]
  \href{https://github.com/talvinckb}{github.com/talvinckb} \enspace·\enspace
  \href{https://www.linkedin.com/in/talvin-ackbaraly}{linkedin.com/in/talvin-ackbaraly}}\\[7pt]
  Final-year engineering student at EPITA, specialized in computer vision and deep learning. I train
  models and ship production software in Python and TypeScript, including a client application I run
  in production as a freelancer. Looking to build LLM-powered products.
\end{minipage}\hfill
\begin{minipage}[c]{\photosize}
  \headerphoto
\end{minipage}

\section{Experience}

\job{Full-Stack Developer \textnormal{\color{soft}(freelance)}}{Sept. 2026 -- ongoing}
    {Self-employed}{\chip{TypeScript}\chip{NestJS}\chip{Next.js}\chip{PostgreSQL}\chip{Sentry}}
    {Inventory management web app for a business client, built end to end and run in production on my
     own: OVH VPS and Object Storage, PostgreSQL backups to Cloudflare R2, Sentry, uptime monitoring.}
\vspace{11pt}
\job{Full-Stack \& DevOps Developer \textnormal{\color{soft}(internship)}}{Sept. 2025 -- Jan. 2026}
    {Studiocanal}{\chip{Java}\chip{Spring Boot}\chip{AWS}\chip{GitLab CI}}
    {Built full-stack features (Spring Boot, Angular) for an internal rights and catalogue platform
     used by 100+ employees. Set up the GitLab CI/CD pipeline of a new module on AWS (EC2, RDS, IAM).}
\vspace{11pt}
\job{IT Consultant \textnormal{\color{soft}(internship)}}{June -- July 2023}
    {Hitachi Solutions France}{\chip{Python}\chip{Django}\chip{UML}}
    {Django prototype of a supplier and contract management app for Sonepar, the firm's largest client.}

\section{Projects}

\proj{Illustration detection for the BnF}{Final-year project, team of 4. Detection and
  hierarchical classification of illustrations in Gallica's documents (IIIF API). In charge of
  detection: training on 10,000 annotated images, benchmarking, hyperparameter tuning,
  \textbf{94.7\%~mAP50}.}
  {\chip{PyTorch}\chip{YOLO}\chip{ConvNeXt}}
\projgap
\proj{Image segmentation with U-Net}{Built from scratch in PyTorch, no pretrained model:
  foreground/background segmentation on COCO, training loop and evaluation on unseen data.}
  {\chip{PyTorch}\chip{U-Net}\chip{COCO}}
\projgap
\proj{Lung function prediction}{Forced vital capacity estimation from CT scans: lung segmentation
  and XGBoost, delivered as a containerized React web app.}
  {\chip{XGBoost}\chip{OpenCV}\chip{Docker}}
\projgap
\proj{License plate recognition}{Classical detection (filters, morphology, HOG) and Random Forest:
  \textbf{86\% precision}. Python vs C++ benchmark.}
  {\chip{OpenCV}\chip{C++}\chip{Random Forest}}
\projgap
\proj{Twitter clone}{Team of 14, from UML design to deployment on Kubernetes: Quarkus back-end,
  React front-end, MongoDB, Neo4j and Elasticsearch.}
  {\chip{Quarkus}\chip{Kubernetes}}


\section{Education}

\edu{EPITA — Engineering School of Computer Science}{Sept. 2022 -- July 2027}
    {Master of Engineering, Image specialization — Paris}
    {Core curriculum: algorithms \& data structures, systems programming (C, Unix), networks, databases.\\
     Image specialization: computer vision, deep learning, GPU computing (CUDA).}
\vspace{7pt}
\edu{Université du Québec à Chicoutimi (UQAC)}{Jan. -- May 2024}
    {Academic exchange semester — Québec, Canada}
    {}


\section{Skills}

\skillrow{Programming}{\chip{Python}\chip{TypeScript}\chip{Java}\chip{C}\chip{C++}\chip{SQL}}
\skillrow{AI \& Machine Learning}{\chip{PyTorch}\chip{Deep learning}\chip{YOLO}\chip{U-Net}\chip{ConvNeXt}\chip{XGBoost}\chip{OpenCV}}
\skillrow{Back-end \& Cloud}{\chip{NestJS}\chip{Spring Boot}\chip{PostgreSQL}\chip{MongoDB}\chip{Elasticsearch}\chip{Docker}\chip{Kubernetes}\chip{AWS}}
\skillrow{AI-assisted development}{\chip{Claude Code}\chip{Codex}\chip{Antigravity}}
\skillrow{Languages}{French (native), English (professional, TOEIC 845)}

\end{document}
